\section{Proofs}\label{app:proofs}

\begin{proof}[Proof of Proposition~\ref{prop:equivalence}]
Suppose substituting the recorded response of $x'$ for $x$ drops nothing owed. By the first condition $x$ denotes one act $a=\llbracket x\rrbracket_t$ and $x'$ one act $a'=\llbracket x'\rrbracket_t$. If $a\ne a'$, the recorded response of $x'$ reflects $a'$ by the second condition, and returning it in place of executing $x$ drops $a$, a contradiction; hence $a=a'$. Conversely, if $a=a'$ the substitution returns the response of the act that $x$ denotes and drops nothing. The test therefore coincides with equality of denotations, which is reflexive, symmetric and transitive.
\end{proof}

\begin{proof}[Proof of Proposition~\ref{prop:lattice}]
If $\mathsf{C}_1,\mathsf{C}_2\subseteq{\approx_t}$ then $\mathsf{C}_1\cup\mathsf{C}_2$ and $\mathsf{C}_1\cap\mathsf{C}_2$ are subsets of $\approx_t$; $\emptyset$ and $\approx_t$ are certifiers and bound the family. The argument for refuters is identical with the complement of $\approx_t$ in place of $\approx_t$. If $\mathsf{C}_1,\dots,\mathsf{C}_k$ are decidable, running their decision procedures and taking the disjunction decides $\bigcup_i\mathsf{C}_i$; decidability need not survive countable unions, so only finite unions are claimed, and only finite unions are used. A relation that is both a certifier and a refuter satisfies $R\subseteq{\approx_t}$ and $R\cap{\approx_t}=\emptyset$, hence $R=\emptyset$.
\end{proof}

\begin{proof}[Proof of Lemma~\ref{lem:scalefree}]
Fix a pool $s$ with distinct entries and suppose $\tau(s)$ is finite and equal to no $s_i$. Then $\tau(s)$ lies in an open interval $I$ that contains no pool value and whose endpoints are consecutive pool values or $\pm\infty$. Let $g$ be the identity outside $I$ and, on $I$, a strictly increasing bijection of $I$ onto itself with $g(\tau(s))\ne\tau(s)$; a piecewise-linear map with one interior breakpoint moved off the diagonal serves when $I$ is bounded, and the dilation $x\mapsto a+2(x-a)$ about its finite endpoint $a$ when it is not. Then $g(s)=s$, so $\tau(g(s))=\tau(s)\ne g(\tau(s))$, contradicting equivariance. Hence $\tau(s)\in\{s_{(1)},\dots,s_{(n)},+\infty,-\infty\}$. For two pools $s,s'$ with distinct entries, the piecewise-linear interpolation of $s_{(j)}\mapsto s'_{(j)}$, extended linearly with slope one beyond the extreme values, is a strictly increasing bijection $g$ of $\mathbb{R}$; $g(s)$ has the same entries as $s'$, so by symmetry $\tau(s')=\tau(g(s))=g(\tau(s))$. If $\tau(s)=s_{(k)}$ then $\tau(s')=s'_{(k)}$, and if $\tau(s)=\pm\infty$ then $\tau(s')=\pm\infty$. The index, or the infinite value, is therefore the same for every pool.
\end{proof}

\begin{proof}[Proof of Theorem~\ref{thm:tight}]
\emph{Reduction.} Let $F$ be the common continuous distribution and put $u_i=F(s_i)$, $u'_i=F(s'_i)$. These are uniform on $(0,1)$ with the same independence structure, and almost surely $s'_i>s_{(k)}$ if and only if $u'_i>u_{(k)}$, because a disagreement requires two sample values in a closed interval on which $F$ is constant, an event of probability zero. We may therefore take $F$ uniform.

\emph{Upper bound.} Condition on the pool. For each $i$ the pool and $u'_i$ are independent, so $P(u'_i>u_{(k)}\mid\text{pool})=1-u_{(k)}=:q$. By the union bound and $P\le1$, $P(\max_iu'_i>u_{(k)}\mid\text{pool})\le\min\{1,mq\}$, whatever the dependence among the $u'_i$. Taking expectations gives $\mathbb{E}[\min\{1,mB\}]$ with $B=1-u_{(k)}\sim\mathrm{Beta}(n+1-k,k)$.

\emph{Attainment.} Let the pool be independent uniforms, let $V$ be uniform and independent of the pool, and put $u'_i=\operatorname{frac}\big(V+(i-1)/m\big)$. Each $u'_i$ is uniform and a function of $V$ alone, so the pool and $u'_i$ are $n+1$ independent uniforms and the law lies in $\mathcal{Q}_{n,m}$. Given the pool, $u'_i>1-q$ if and only if $V$ lies in an arc of length $q$ of the circle $\mathbb{R}/\mathbb{Z}$ whose right endpoint is $1-(i-1)/m$. The $m$ right endpoints are equally spaced at distance $1/m$. If $q\le1/m$ the arcs are pairwise disjoint and their union has measure $mq$; if $q>1/m$ consecutive arcs overlap and the union is the whole circle. Hence $P(\max_iu'_i>u_{(k)}\mid\text{pool})=\min\{1,mq\}$, and the upper bound holds with equality.

\emph{The fence's threshold.} With $k=k^\ast$, $B\sim\mathrm{Beta}(\rho,n+1-\rho)$ and $\mathbb{E}[B]=\rho/(n+1)$. Since $\min\{1,mB\}=mB-(mB-1)_+$, the worst case is $m\rho/(n+1)-\mathbb{E}[(mB-1)_+]$, and $m\rho/(n+1)\le\alpha$ by Proposition~\ref{prop:dilemma}(b). Because $B\le1$, $(mB-1)_+\le(m-1)\,\mathbf{1}\{B>1/m\}$. The variable $B$ is distributed as the $\rho$-th smallest of $n$ independent uniforms, so $B>1/m$ exactly when fewer than $\rho$ of them fall at or below $1/m$: $P(B>1/m)=P(X\le\rho-1)$ with $X\sim\mathrm{Bin}(n,1/m)$. Put $\mu=n/m$. Then $\rho-1\le(n+1)\alpha/m-1<\alpha\mu$, and the Chernoff bound $P(X\le(1-\delta)\mu)\le e^{-\delta^2\mu/2}$ with $\delta=1-\alpha$ gives $P(B>1/m)\le\exp(-(1-\alpha)^2n/(2m))$.
\end{proof}

\begin{proof}[Proof of Corollary~\ref{cor:tight}]
(a) For $m=1$, $(B-1)_+=0$, so the worst case is $\rho/(n+1)$, which is $\alpha$ when $(n+1)\alpha$ is an integer. For general $m$, $m\rho/(n+1)\ge\alpha-m/(n+1)$ and the correction term of Theorem~\ref{thm:tight} tends to zero, so the worst case tends to $\alpha$ and exceeds any $\alpha'<\alpha$ for large $n$.

(b) For $m=1$ the worst case of rank $k$ is $\mathbb{E}[B]=(n+1-k)/(n+1)$, and $k<k^\ast$ gives $n+1-k\ge\rho+1>(n+1)\alpha$. For a fixed $\gamma\in(\alpha/m,1)$ let $k_\gamma=\lceil(n+1)(1-\gamma)\rceil$, which lies in $\{1,\dots,n\}$ for large $n$, and let $B_\gamma\sim\mathrm{Beta}(n+1-k_\gamma,k_\gamma)$. Its mean tends to $\gamma$ and its variance is at most $1/(4(n+2))$, so $B_\gamma\to\gamma$ in probability; as $x\mapsto\min\{1,mx\}$ is bounded and continuous, the worst case tends to $\min\{1,m\gamma\}$, which exceeds $\alpha$ because $m\gamma>\alpha$ and $\alpha<1$. The level $\alpha/m$ itself is valid for every $n$ by Theorem~\ref{thm:bound}.

(c) Let $m>\alpha(n+1)$ and $k\le n$. On the same pool $u_{(k)}\le u_{(n)}$, so $1-u_{(k)}\ge1-u_{(n)}$ and the worst case of rank $k$ is at least that of rank $n$, for which $B\sim\mathrm{Beta}(1,n)$ and $P(B>b)=(1-b)^n$. Then
\begin{align*}
\mathbb{E}[\min\{1,mB\}]&=\int_0^1P(mB>u)\,du=\int_0^1\Big(1-\frac um\Big)^n du\\
&=\frac{m}{n+1}\Big(1-\Big(1-\frac1m\Big)^{n+1}\Big).
\end{align*}
Put $x=m/(n+1)$. Since $(1-1/m)^{n+1}\le e^{-(n+1)/m}=e^{-1/x}$, this is at least $h(x)=x(1-e^{-1/x})$. The function $h$ is continuous and strictly increasing on $(0,\infty)$, because $h'(x)=1-e^{-1/x}(1+1/x)>0$; it tends to $1$ as $x\to\infty$; and $h(\alpha)<\alpha$. Hence there is a unique $x_\alpha>\alpha$ with $h(x_\alpha)=\alpha$. With $\delta_\alpha=x_\alpha-\alpha$, every $x>\alpha+\delta_\alpha$ gives a worst case above $\alpha$. Solving numerically, $\delta_{0.05}=1.03\times10^{-10}$ and $\delta_{0.01}=1.7\times10^{-18}$. For $\alpha=1/20$, $m>(n+1)/20$ implies $20m\ge n+2$, so $x-\alpha\ge1/(20(n+1))$, which exceeds $\delta_{0.05}$ whenever $n+1<4.8\times10^8$. The threshold $-\infty$ suppresses every request. Under the stated condition, every scale-free threshold other than $+\infty$ is therefore invalid on $\mathcal{Q}_{n,m}$.
\end{proof}

\begin{proof}[Proof of Corollary~\ref{cor:mondrian}]
(a) Fix a class $c$ and condition on it. The proof of Theorem~\ref{thm:bound} uses only that each distinct comparison score of the request is exchangeable with the pool on which the threshold is computed; under A6 within class $c$ this holds for the class-$c$ pool of size $n_c$, and the argument goes through with $n_c$ in place of $n$. (b) The first statement is Proposition~\ref{prop:dilemma}(a) with $n_c$ in place of $n$. Every class is certifiable iff each $n_c\ge\lceil m/\alpha\rceil-1$, i.e.\ iff $\min_c n_c$ is; summing over the $K$ classes gives $n\ge K(\lceil m/\alpha\rceil-1)$. Since $\min_c n_c\le n/K$ and $x\mapsto\lfloor\alpha(x+1)\rfloor$ is non-decreasing, $\min_c\lfloor\alpha(n_c+1)\rfloor\le\lfloor\alpha(n/K+1)\rfloor$. (c) Let the request be of class $c$. If $S^\ast>\max_{c'}t_{c'}(m)$ then $S^\ast>t_c(m)$, so the suppression event under the maximum is contained in the suppression event under $t_c(m)$, whose probability given class $c$ is at most $\alpha$ by (a). If some class is not certifiable, $\max_{c'}t_{c'}(m)=+\infty$.
\end{proof}

\begin{proof}[Proof of Proposition~\ref{prop:impossible}]
Let a threshold rule suppress iff the score exceeds $\tau$. Zero wrong suppression requires that $F_{\mathrm{diff}}$ place no mass above $\tau$. Zero silent escape without a band requires that $F_{\mathrm{same}}$ place no mass at or below $\tau$. Together these require that every score with positive probability under $F_{\mathrm{diff}}$ be at most $\tau$ and every score with positive probability under $F_{\mathrm{same}}$ exceed $\tau$, contradicting overlap. With a band $(\tlo,\thi]$, silent escape can be removed by widening the band over the overlap, but the duplicates in the overlap are then flagged rather than suppressed, so deduplication falls by the overlap mass; the error is relocated, not removed. A decision that does not consult the score is not subject to this argument.
\end{proof}

\begin{proof}[Proof of Proposition~\ref{prop:amo}]
Consider the first execution for $\mathit{op}$ to reach the insert. By A3 the insert either creates the commit row or observes an existing one. Any later execution computing $\mathit{op}$---in the same process or a fresh one---computes the same identifier by A1 and A2, finds the committed row, and returns the recorded response on the certifier branch without dispatching. By A7 no second execution is inside the interval between its lookup and its insert while the first is, so the dispatch branch is entered at most once before a commit row exists, and by A4 that branch dispatches the effect once and commits only on success. If the process dies after the dispatch and before the insert, no commit row exists, the intent remains \textsc{pending}, and a later execution finds no row and may dispatch again; this is the stated exception.
\end{proof}

\begin{proof}[Proof of Proposition~\ref{prop:compose}]
\emph{Case A.} If the fence suppresses, nothing is dispatched. If it executes and the process dies inside the commit window, a later execution re-dispatches the same act, possibly paraphrased. By paraphrase invariance both dispatches carry the same key, so by A8 the provider applies at most one effect; by completeness no distinct act carries that key, so no legitimate act is absorbed. Outside the window, Proposition~\ref{prop:amo} applies. If the schema satisfies A5$''$, $\sigma_t$ is both invariant and complete, so $\pi_t=\sigma_t$ qualifies.
\emph{Case B.} Without A5$''$ there exist $x\not\approx_t x'$ with $\sigma_t(x)=\sigma_t(x')$; presented to a provider obeying A8 under key $\sigma_t$, the second is absorbed although it is a distinct act, which is a wrong suppression invisible to the local log. The canonical identifier is admissible by Proposition~\ref{prop:sound} but changes under paraphrase, because $\kappa_t$ normalises free text without interpreting it, so a paraphrased retry presents a new key and is applied.
\end{proof}

\begin{proof}[Proof of Proposition~\ref{prop:structure}]
If $x\approx_t x'$ but the fence computes different identifiers, the log's uniqueness constraint does not fire and a provider keyed on those identifiers receives two keys; under A8 it is obliged to apply both. If $x\not\approx_t x'$ but they receive the same identifier, both the log and the provider absorb the second, and both behave as specified. In either case the failure is attributable to identity recognition alone. Conversely, under Case A of Proposition~\ref{prop:compose} the provider applies at most one effect per act whatever the local fence serialised, so A7 is not needed for at-most-once application.
\end{proof}
